\SourceLeaf{8}
quam difficillimam problematis partem absolvere contigit ill. \textsc{Legendre} (\textit{Recherches sur l’attraction des sphéroïdes homogènes, Mémoires présentés à l’acad. roi. des sc. T.} X.)

Disquisitionem generalissimam de attractione sphaeroidum non per revolutionem ortarum, sed quarum sectiones cum quolibet plano sunt ellipses, iamiam inchoaverat \textsc{Mac Laurin}, sed substiterat in attractione punctorum in aliquo trium axium positorum. Theorema principale, cui solutio problematis generalissima praesertim innititur, per inductionem quidem iam coniectaverat ill. \textsc{Legendre} in commentatione modo laudata, sed ill. \textsc{Laplace} primo successit, omnia rigorose demonstrare atque sic solutionem ab omni parte perfectam reddere (\textit{Hist. de l’acad. roi. des sc. de Paris} 1782; eadem solutio repetita in operibus \textit{Théorie du mouvement et de la figure elliptique des planètes}, atque \textit{Mécanique céleste} Vol. 2).

\expandafter\def\csname HistoricalNoteText008\endcsname{De his duabus solutionibus e. g. ita iudicat ill. \textsc{Lagrange}: \textit{On ne peut regarder leurs solutions que comme des chefs-d’œuvres d’analyse, mais on peut désirer encore une solution plus directe et plus simple; et les progrès continuels de l’analyse donnent lieu de l’espérer.} Nouv. Mém. de l’acad. de Berlin 1793. p. 263.}
Elegantiam ingeniique subtilitatem in hac ill. \textsc{Laplace} solutione eminentem nemo quidem non mirabitur: nihilominus tamen ipsa subtilitas arsque admiranda, per quam arduas difficultates superavit, geometris desiderium liquit solutionis simplicioris, minus intricatae magisque directae. Nec plane satisfecit huic desiderio ill. \textsc{Legendre} per novam theorematis principalis demonstrationem (\textit{Hist. de l’acad. roi. des sc.} 1788, \textit{Sur les intégrales doubles}), etiamsi exquisita ars analytica omnium geometrarum suffragia merito tulerit\HistoricalNote{008}. Postea clar. \textsc{Biot} solutionem alteram, alteram clar. \textsc{Plana} simpliciorem reddere conati sunt (\textit{Mém. de l’institut T.} VI, \textit{Memorie di matematica e di fisica della società italiana T.} XV): sed sic quoque utramque solutionem ad intricatissimas analyseos applicationes referendam esse, quisque facile concedet.

Gratam itaque analystis atque astronomis fore speramus solutionem novam problematis celebratissimi per viam plane diversam procedentem, et ni fallimur ea simplicitate gaudentem, ut nihil amplius desiderandum linquat.

Ipsa quidem solutio nostra paucissimis pagellis continebitur. Operae tamen pretium esse censemus, antequam ad ipsum problema, cui haec commentatio dicata est, descendamus, quasdam disquisitiones praeliminares, quae in aliis quoque occasionibus opportune applicari poterunt, aliquanto generalius exsequi, fusiusque explicare, quam instituti nostri ratio per se spectata postularet.
